\chapter{Qué distingue a la dimensión finita}\label{chap:22}

En $\R^n$ estamos acostumbrados a que todas las normas describan la misma convergencia, a que toda transformación lineal sea continua y a que los conjuntos cerrados acotados sean compactos. Estas propiedades no son consecuencias exclusivas de la norma: dependen de la dimensión finita. El objetivo de este capítulo es identificar dónde interviene esa hipótesis y qué propiedades permiten reconocerla.

Usaremos el lema de coordenadas del capítulo~\ref{chap:19}, que ya se demostró mediante la compacidad de la esfera de $\F^n$, la caracterización de funcionales continuas por su núcleo y el lema de Riesz. No se presupone ningún teorema de aplicación abierta para espacios de Banach.

\section{Subespacios finitos y continuidad}

\begin{proposicion}\label{ch22:suma-cerrada}
Si $M$ es un subespacio cerrado de un espacio normado $V$ y $F$ es un subespacio de dimensión finita, entonces $M+F$ es cerrado.
\end{proposicion}
\begin{proof}
Empecemos por $F=\F u$. Si $u\in M$, la suma es $M$ y no hay nada que probar. Si $u\notin M$, sea $d=\dist(u,M)>0$. Para cualquier escalar $t$, la invariancia de $M$ bajo multiplicación da
\[
 \dist(tu,M)=|t|d;
\]
para $t=0$ ambos lados son cero, y para $t\ne0$ se obtiene cambiando $m$ por $m/t$ en el ínfimo.

Sea $z_n=m_n+t_nu$ una sucesión de $M+\F u$ que converge a $z$ en $V$. Como $m_n-m_k\in M$,
\[
 |t_n-t_k|d\leq\norm{z_n-z_k}.
\]
La convergencia de $(z_n)$ implica que es de Cauchy, de modo que $(t_n)$ es de Cauchy en $\F$ y converge a algún $t$. Entonces $m_n=z_n-t_nu\to z-tu$. Por cerradura de $M$, $z-tu\in M$, y así $z\in M+\F u$. Esto prueba que la suma es cerrada.

Para $F=\spanop\{u_1,\ldots,u_k\}$, aplique el caso anterior sucesivamente a
\[
 M,\quad M+\F u_1,\quad M+\F u_1+\F u_2,\quad\ldots.
\]
Cada paso añade una dirección a un subespacio ya cerrado. Después de un número finito de pasos se obtiene $M+F$ cerrado.
\end{proof}

\begin{corolario}\label{ch22:finito-cerrado}
Todo subespacio de dimensión finita de un espacio normado es cerrado.
\end{corolario}
\begin{proof}
Tome $M=\{0\}$ en la proposición anterior. Esta conclusión coincide con la obtenida por coordenadas en el lema~\ref{ch19:coordenadas}; la nueva prueba resalta la estabilidad de la cerradura al añadir un número finito de direcciones.
\end{proof}

\begin{teorema}\label{ch22:rango-finito-kernel}
Sea $T:V\to W$ lineal, con $T(V)$ de dimensión finita. Entonces $T$ es continuo si y solo si $\ker T$ es cerrado.
\end{teorema}
\begin{proof}
La necesidad se sigue de que $\ker T=T^{-1}(\{0\})$. Para la suficiencia, si $T(V)=\{0\}$, $T$ es el operador cero. En otro caso, elija una base $w_1,\ldots,w_k$ de $T(V)$ y vectores $v_j\in V$ con $Tv_j=w_j$. Para cada $v$ existen coordenadas únicas $\varphi_j(v)$ tales que
\[
 Tv=\sum_{j=1}^k\varphi_j(v)w_j.
\]
La unicidad de coordenadas y la linealidad de $T$ implican que cada $\varphi_j$ es lineal. Afirmamos
\[
 \ker\varphi_i=\ker T+\spanop\{v_j:j\ne i\}.
\]
Si $\varphi_i(v)=0$, el vector $v-\sum_{j\ne i}\varphi_j(v)v_j$ tiene imagen cero por $T$, de modo que pertenece a $\ker T$. Esto prueba una inclusión. En la otra dirección, la imagen de $m+\sum_{j\ne i}a_jv_j$, con $m\in\ker T$, es $\sum_{j\ne i}a_jw_j$, cuya coordenada $i$ es cero.

La proposición~\ref{ch22:suma-cerrada} muestra que $\ker\varphi_i$ es cerrado. Por la proposición~\ref{ch21:kernel-cerrado}, $\varphi_i$ es continua. Finalmente,
\[
 \norm{Tv}\leq\sum_{j=1}^k|\varphi_j(v)|\norm{w_j}
 \leq\left(\sum_{j=1}^k\norm{\varphi_j}\norm{w_j}\right)\norm v.
\]
Esta cota demuestra continuidad de $T$.
\end{proof}

\begin{corolario}\label{ch22:lineal-finito}
Si $V$ tiene dimensión finita, toda transformación lineal $T:V\to W$ es continua. Todas las normas en $V$ son equivalentes y $V$ es de Banach.
\end{corolario}
\index{dimension finita@dimensión finita!equivalencia de normas}
\begin{proof}
Si $V=\{0\}$, todas las conclusiones son inmediatas. Supongamos $\dim V>0$. El núcleo de $T$ es un subespacio finito dimensional y, por ello, cerrado. Su imagen también tiene dimensión finita, como se ve tomando las imágenes de una base de $V$. El teorema anterior implica que $T$ es continua.

Si $\norm{\cdot}_1$ y $\norm{\cdot}_2$ son dos normas en $V$, la identidad de $(V,\norm{\cdot}_1)$ en $(V,\norm{\cdot}_2)$ y la identidad inversa son lineales. Por lo recién probado, ambas son continuas. Existen por ello $c,C>0$ tales que $c\norm v_1\leq\norm v_2\leq C\norm v_1$. Estas desigualdades son la equivalencia de normas.

Elija una base de $V$ y la norma de coordenadas $\norm{\sum a_jv_j}_0=\max_j|a_j|$. Una sucesión de Cauchy para esa norma es de Cauchy en cada coordenada; la completitud de $\F$ produce sus límites y, por tanto, un límite en $V$. Así $\norm{\cdot}_0$ es completa. Las desigualdades de equivalencia muestran que una sucesión de Cauchy en la norma original lo es en $\norm{\cdot}_0$, y que la convergencia en esta implica convergencia en la original. Por tanto $V$ es de Banach.
\end{proof}

\begin{proposicion}[Estabilidad de la independencia]\label{ch22:independencia}
Si $v_1,\ldots,v_k$ son linealmente independientes en $V$, existe $\delta>0$ tal que $\norm{w_j-v_j}<\delta$ para todo $j$ implica que $w_1,\ldots,w_k$ también son linealmente independientes.
\end{proposicion}
\begin{proof}
El lema~\ref{ch19:coordenadas} proporciona $c>0$ con
\[
 \norm{\sum_{j=1}^k a_jv_j}\geq c\max_j|a_j|
 \geq\frac ck\sum_{j=1}^k|a_j|.
\]
Tome $\delta=c/(2k)$. Si $\sum_j a_jw_j=0$, entonces
\begin{align*}
 \frac ck\sum_j|a_j|
 &\leq\norm{\sum_j a_jv_j}
 =\norm{\sum_j a_j(v_j-w_j)}\\
 &\leq\sum_j|a_j|\norm{v_j-w_j}
 \leq\frac c{2k}\sum_j|a_j|.
\end{align*}
Esto fuerza $\sum_j|a_j|=0$, y por ello todos los coeficientes son cero.
\end{proof}

\section{Compacidad y dimensión}

\begin{teorema}[Heine--Borel en dimensión finita]\label{ch22:heine-borel}
En un espacio normado de dimensión finita, un conjunto es compacto si y solo si es cerrado y acotado.
\end{teorema}
\begin{proof}
En todo espacio métrico, los compactos son cerrados y acotados. Para el recíproco, use una base para identificar $V$ linealmente con $\F^k$. El corolario~\ref{ch22:lineal-finito} muestra que la identificación y su inversa son continuas y que las normas se acotan mutuamente. Por tanto la imagen de un cerrado acotado de $V$ es cerrada y acotada en $\F^k$. Heine--Borel euclidiano la hace compacta, y su imagen por la identificación inversa es compacta en $V$.
\end{proof}

\begin{teorema}\label{ch22:caracterizacion}
Para un espacio normado $V$, las siguientes condiciones son equivalentes:
\begin{enumerate}[label=(\roman*)]
\item $V$ tiene dimensión finita;
\item la bola cerrada $\overline B(0,1)$ es compacta;
\item $V$ es localmente compacto.
\end{enumerate}
\end{teorema}
\index{compacidad!local en espacios normados}
\begin{proof}
Si la dimensión es finita, Heine--Borel da la compacidad de todas las bolas cerradas. Cada punto tiene una de esas bolas como vecindad compacta, por lo que (i) implica (ii) y (iii).

Si $V$ es localmente compacto, el origen tiene una vecindad compacta $K$. Existe $r>0$ con $B(0,r)\subseteq K$. La bola cerrada $\overline B(0,r/2)$ es un cerrado de $V$ contenido en $K$, luego es compacta. La homotecia $v\mapsto (2/r)v$ muestra que $\overline B(0,1)$ es compacta. Esto prueba (iii)$\Rightarrow$(ii).

Supongamos (ii). La compacidad implica acotación total, así que la bola unitaria está cubierta por un número finito de bolas abiertas de radio $1/2$, con centros $v_1,\ldots,v_k\in V$. Sea $M=\spanop\{v_1,\ldots,v_k\}$. Es cerrado por ser de dimensión finita. Si $M\ne V$, el lema de Riesz daría $v$ con $\norm v=1$ y $\dist(v,M)>1/2$. Sin embargo, algún centro satisface $\norm{v-v_j}<1/2$, contradicción. Por tanto $M=V$ y $V$ tiene dimensión finita.
\end{proof}

\begin{corolario}\label{ch22:no-sigma-compacto}
Un espacio de Banach infinito dimensional no es unión numerable de conjuntos compactos.
\end{corolario}
\begin{proof}
Si $V=\bigcup_nK_n$ con cada $K_n$ compacto, todos son cerrados. Baire da un $K_n$ con interior no vacío. Existe una bola $B(x,r)\subseteq K_n$. La bola cerrada $\overline B(x,r/2)$ es un cerrado contenido en el compacto $K_n$, y por ello es compacta. Trasladar y reescalar hace compacta la bola unitaria de $V$. El teorema anterior contradice la dimensión infinita.
\end{proof}

\begin{ejemplo}
En $\ell^2$, los vectores canónicos $e_n$ tienen norma $1$ y $\norm{e_n-e_m}=\sqrt2$ si $n\ne m$. La sucesión no tiene ninguna subsucesión de Cauchy y, por tanto, ninguna convergente. La bola unitaria es cerrada y acotada, pero no compacta. El teorema anterior demuestra que siempre existe una obstrucción de este tipo en dimensión infinita, aunque no haya una base canónica disponible.
\end{ejemplo}

\section{Ejercicios resueltos}

\begin{ejercicio}[Lectura 22, ejercicio 1]
Demuestre que toda transformación lineal sobreyectiva de un espacio normado sobre un espacio normado de dimensión finita es abierta.
\end{ejercicio}
\begin{solucion}
Sea $T:V\to W$ sobreyectiva. Si $W=\{0\}$, la imagen de un abierto es $\varnothing$ o $\{0\}$, ambos abiertos en $W$. Supongamos $\dim W=k>0$. Elija una base $w_1,\ldots,w_k$ y preimágenes $v_j$ con $Tv_j=w_j$. Defina
\[
 R\left(\sum_j a_jw_j\right)=\sum_j a_jv_j.
\]
Es una aplicación lineal con $TR=I_W$. Como su dominio es finito dimensional, es continua; sea $C=\norm R>0$. Si $U\subseteq V$ es abierto y $y=Tx\in T(U)$ con $x\in U$, existe $r>0$ con $B(x,r)\subseteq U$. Para $z\in W$ con $\norm{z-y}<r/C$,
\[
 \norm{R(z-y)}\leq C\norm{z-y}<r.
\]
Por ello $x+R(z-y)\in U$, y su imagen es $y+(z-y)=z$. Así $B(y,r/C)\subseteq T(U)$ y $T(U)$ es abierto.

No se supuso que $T$ fuera continua. La apertura procede de la existencia de una sección lineal continua $R$ con dominio finito dimensional; no es una aplicación del teorema de la aplicación abierta para Banach.
\end{solucion}

\begin{ejercicio}[Lectura 22, ejercicio 2]
Sea $V$ finito dimensional con base $v_1,\ldots,v_k$. Pruebe que existe $C>0$ tal que
\[
 \norm T\leq C\max_j\norm{Tv_j}
\]
para todo operador lineal $T$ de $V$. Deduzca que si $T_n$ converge puntualmente a cero, entonces $\norm{T_n}\to0$.
\end{ejercicio}
\begin{solucion}
Por el lema de coordenadas existe $c>0$ tal que $\norm{\sum_j a_jv_j}\geq c\max_j|a_j|$. Para $v=\sum_j a_jv_j$,
\[
 \norm{Tv}\leq\sum_j|a_j|\norm{Tv_j}
 \leq\frac kc\norm v\max_j\norm{Tv_j}.
\]
Tomando el supremo sobre $\norm v\leq1$, sirve $C=k/c$. El argumento funciona también si el codominio de $T$ es cualquier otro espacio normado. Si $T_n$ converge puntualmente a cero, cada una de las $k$ sucesiones $T_nv_j$ converge a cero. El máximo de sus normas tiende a cero porque el conjunto de índices es finito. La desigualdad recién probada concluye que $\norm{T_n}\to0$. Para $V=\{0\}$ todos los operadores son cero y la conclusión es directa.
\end{solucion}

\begin{ejercicio}[Lectura 22, ejercicio 3]
Sea $\mathcal A_n=\{T\in\Bcal(V,W):\dim T(V)\leq n\}$. Demuestre que $\mathcal A_n$ es cerrado en norma de operador.
\end{ejercicio}
\begin{solucion}
Suponga $T_j\in\mathcal A_n$ y $T_j\to T$ en norma. Si $\dim T(V)>n$, existen $v_1,\ldots,v_{n+1}$ tales que $Tv_1,\ldots,Tv_{n+1}$ son linealmente independientes. Por la proposición~\ref{ch22:independencia}, existe $\delta>0$ que conserva esa independencia bajo perturbaciones menores que $\delta$. Para cada $i$,
\[
 \norm{T_jv_i-Tv_i}\leq\norm{T_j-T}\norm{v_i}\to0.
\]
Hay un mismo $j$ suficientemente grande para que las $n+1$ diferencias sean menores que $\delta$. Entonces $T_jv_1,\ldots,T_jv_{n+1}$ son independientes, contradiciendo $\dim T_j(V)\leq n$. Así $T\in\mathcal A_n$. El criterio de sucesiones en el espacio normado $\Bcal(V,W)$ prueba la cerradura.
\end{solucion}

\begin{ejercicio}[Lectura 22, ejercicio 4]
Si $M$ tiene dimensión finita y $v\notin M$, demuestre que existe $m_0\in M$ con $\dist(v,M)=\norm{v-m_0}$. Deduzca que existe $u$ con $\norm u=\dist(u,M)=1$.
\end{ejercicio}
\begin{solucion}
Sea $d=\dist(v,M)$. El subespacio es cerrado, por lo que $d>0$. Elija $m_j\in M$ con $\norm{v-m_j}\to d$. Como $0\in M$, $d\leq\norm v$. Podemos descartar términos iniciales para que $\norm{v-m_j}\leq\norm v+1$. Entonces
\[
 \norm{m_j}\leq\norm v+\norm{v-m_j}\leq2\norm v+1.
\]
La bola correspondiente de $M$ es compacta por Heine--Borel. Una subsucesión converge a $m_0\in M$, y la continuidad de la norma da $\norm{v-m_0}=d$.

Defina $u=(v-m_0)/d$. Su norma es $1$. Como trasladar por un elemento de $M$ no cambia la distancia a $M$, y multiplicar por $1/d$ la multiplica por $1/d$,
\[
 \dist(u,M)=\frac{\dist(v-m_0,M)}d=\frac{\dist(v,M)}d=1.
\]
En dimensión finita la distancia se alcanza. El lema de Riesz general solo garantiza distancias arbitrariamente próximas a $1$, porque un subespacio cerrado infinito dimensional puede no tener mejores aproximantes.
\end{solucion}

\begin{ejercicio}[Lectura 22, ejercicio 5]
Sea $c_{00}=\spanop\{e_n:n\in\N\}$, el espacio de sucesiones de soporte finito. Demuestre que no es completo con ninguna norma.
\end{ejercicio}
\begin{solucion}
Fije una norma arbitraria y sea $M_n=\spanop\{e_1,\ldots,e_n\}$. Cada $M_n$ es cerrado porque tiene dimensión finita. Es propio y, por el ejercicio 3 del capítulo~\ref{chap:19}, tiene interior vacío. Todo vector de $c_{00}$ tiene soporte finito, de modo que $c_{00}=\bigcup_{n=1}^{\infty}M_n$. Si la norma elegida fuera completa, esta unión contradiría Baire. Como no se impuso ninguna condición especial a la norma, ninguna puede hacer completo a $c_{00}$.
\end{solucion}

\begin{ejercicio}[Lectura 22, ejercicio 6]
Sea $V\subseteq\ell^1$ de dimensión finita. Demuestre que para todo $\eps>0$ existe $N\in\N$ tal que, para todo $x\in V$ y todo $\lambda\in\F$,
\[
 \norm{x+\lambda e_N}_1\geq|\lambda|+(1-\eps)\norm x_1.
\]
\end{ejercicio}
\begin{solucion}
Para cada $n$ considere la funcional de coordenada $\varphi_n(x)=x_n$ restringida a $V$. Es continua y $|x_n|\leq\norm x_1$. Como toda sucesión absolutamente sumable tiene términos que tienden a cero, $\varphi_n(x)\to0$ para cada $x\in V$. La versión del ejercicio 2 con codominio $\F$ muestra que $\norm{\varphi_n}\to0$. Elija $N$ con $\norm{\varphi_N}\leq\eps/2$. Entonces $|x_N|\leq(\eps/2)\norm x_1$ para todos los $x\in V$ a la vez. Finalmente,
\begin{align*}
 \norm{x+\lambda e_N}_1
 &=\sum_{n\ne N}|x_n|+|x_N+\lambda|\\
 &\geq\norm x_1-|x_N|+|\lambda|-|x_N|\\
 &\geq|\lambda|+(1-\eps)\norm x_1.
\end{align*}
El índice $N$ depende de $V$ y $\eps$, pero no de $x$ ni de $\lambda$. Esa uniformidad es exactamente lo que aporta la dimensión finita.
\end{solucion}

\section{Perspectiva}
En dimensión finita, un número finito de coordenadas controla simultáneamente todas las direcciones. De ahí proceden la equivalencia de normas, la continuidad automática y la compacidad de la bola. En dimensión infinita ya no puede sustituirse una bola por un número finito de direcciones con un error uniformemente pequeño; el lema de Riesz hace precisa esa obstrucción.
